

%\section{\uppercase{Conclusion}}

%Our study explores the double descent phenomenon in three recent deep learning models. QNNs exhibited clear double descent on CIFAR-10, while LTCs only showed some evidence of it, requiring further validation. SNNs did not display double descent. We found that learning rate schedulers, optimizers, and training epochs significantly influence double descent. Future work will involve broader experiments with varied parameters, optimizers, schedulers, and models, including FNNs, RNNs, and GANs.

%We believe that double descent reflects the behavior of dynamically learning features by the model during training, signifying that the model first learns shallow features and subsequently captures deeper, more complex features. Models exhibiting this pattern often demonstrate strong generalization. However, the absence of the double descent phenomenon does not necessarily indicate good or poor generalization performance. For instance, a model may quickly learn sufficient features, causing the test error curve to decline rapidly and stabilize at a low value without exhibiting a second increase. Conversely, the test error curve may take on a U-shaped pattern, as observed in SNN models trained without a learning rate scheduler. The reasons behind the lack of double descent phenomenon in SNNs, however, require further investigation.

%In summary, we believe that double descent is generally a favorable phenomenon for generalization. However, researchers should not explicitly aim to achieve double descent; instead, they should remain focused on the ultimate goal of machine learning—achieving better generalization. To this end, selecting appropriate hyperparameters and adopting efficient learning rate schedulers, as demonstrated in this study, can be effective strategies.


This research explored the performance of Bayesian Neural Networks (BNNs) robustness against Single Event Upsets (SEUs). Our initial findings show that the Deterministic CNN model is better in test accuracy compared to the Bayesian CNN model. On the other hand, Bayesian CNN model is more robust, looking at its performance in ARIn and Softmax Difference. Based on our experiment result, model with dropout is the most robust for Bayesian CNN. The Gaussian prior is the most robust. SEU against the prior center parameter is more robust than the spread parameter. Sigmoid offers better robustness in ASD and ARIn, RWG is the most robust in AAAD. Trade-off between test accuracy and robustness is observed. Convolutional layers are more robust than the fully connected layer. Attack on the first exponent bit is the most devastating; mantissa bit flips had a minimum effect. Flipping a bit from 1 to 0 is found to be far more robust.

%Future research could move beyond evaluating performance on unseen datasets (e.g., accuracy, F1 score) and instead analyze how SEUs affect the magnitudes and distributions of weights and biases in subsequent layers. Such an analysis may help identify sensitive components within the network and inform targeted mitigation strategies.

%Beyond architectural modifications, future work could also investigate distributional adjustments at the predictive level to improve robustness. One promising direction is applying Laplace smoothing to the output probability distributions to stabilize predictions under SEU-induced perturbations. It would be interesting if we can find any correlation between the smoothing factor and SEU robustness.